\documentclass[12pt,a4paper]{article}
\usepackage[utf8]{inputenc}
\usepackage[T2A]{fontenc}
\usepackage[russian]{babel}
\usepackage{amsmath, amssymb, mathtools}
\usepackage{array, booktabs}
\usepackage[hidelinks]{hyperref}
\usepackage{geometry}
\geometry{margin=2cm}
\pagestyle{empty}

\begin{document}

% ─────────────────────────────────────────────────────────────────────────────
\begin{center}
{\large\bfseries Слайд 1. Паспорт источника и формирование BPA}
\end{center}
\medskip
\hrule
\bigskip

\noindent
\textbf{Рамка рассуждений:}
\[
  \Theta = \{T,\,F\}, \qquad
  2^\Theta = \bigl\{\emptyset,\;\{T\},\;\{F\},\;\Theta\bigr\}
\]

\noindent
\textbf{Множество источников:}
$\mathcal{S} = \{\mathrm{vac},\;\mathrm{exp},\;\mathrm{fc}\}$.
Каждый источник $s \in \mathcal{S}$ задаётся кортежем
$(n_s^{\mathrm{tot}},\; n_s^{+},\; \bar{a}_s^{+},\; n_s^{-},\;
 \bar{a}_s^{-},\; \alpha_s,\; w_s,\; \lambda_s)$.

\bigskip
\noindent\textbf{Шаг 1. Коэффициент уверенности $\kappa_s$}

\[
  n_s^{\mathrm{eff}} = n_s^{\mathrm{rel}} \cdot \alpha_s,
  \qquad
  \kappa_s = \frac{n_s^{\mathrm{eff}}}{n_s^{\mathrm{eff}} + \lambda_s}
  \;\in (0,1)
\]

\bigskip
\noindent\textbf{Шаг 2. Базовое распределение вероятности (BPA)}\\
Единая схема для \emph{всех} источников:

\begin{align*}
  m_s\!\bigl(\{T\}\bigr)
    &= \kappa_s \cdot \bar{a}_s^{+} \cdot
       \frac{n_s^{+}}{n_s^{\mathrm{tot}}}
    &&\text{(масса «важно»)}\\[6pt]
  m_s\!\bigl(\{F\}\bigr)
    &= \kappa_s \cdot \bar{a}_s^{-} \cdot
       \frac{n_s^{-}}{n_s^{\mathrm{tot}}}
    &&\text{(масса «не важно»)}\\[6pt]
  m_s(\Theta)
    &= 1 - m_s\!\bigl(\{T\}\bigr) - m_s\!\bigl(\{F\}\bigr)
    &&\text{(неопределённость)}
\end{align*}

\noindent
Для \textbf{vac, exp}: нет отрицательных сигналов $\Rightarrow$
$n_s^{-} = 0 \Rightarrow m_s(\{F\}) = 0$ (следствие данных, не аксиома).\\
Для \textbf{fc} (смешанный): $n_{\mathrm{fc}}^{+} > 0$ и
$n_{\mathrm{fc}}^{-} > 0$ $\Rightarrow$ оба слагаемых ненулевые.

\bigskip
\noindent\textbf{Шаг 3. Дисконтирование по надёжности $w_s$}

\begin{align*}
  m'_s(A)      &= w_s \cdot m_s(A),
    \quad A \neq \Theta \\[6pt]
  m'_s(\Theta) &= (1 - w_s) + w_s \cdot m_s(\Theta)
\end{align*}

\smallskip
\begin{center}
\begin{tabular}{lcc}
\toprule
Источник & $w_s$ & $\lambda_s$ \\
\midrule
Эксперты & $0{,}90$ & $5$ \\
Вакансии & $0{,}80$ & $15$ \\
Прогнозы & $0{,}60$ & $2$ \\
\bottomrule
\end{tabular}
\end{center}

\medskip
\noindent\textit{Выход: дисконтированные тройки
$(m'_s(\{T\}),\; m'_s(\Theta),\; m'_s(\{F\}))$
для каждого $s \in \mathcal{S}$; сумма равна 1.
Передаются на комбинирование по правилу Демпстера.}

\newpage
% ─────────────────────────────────────────────────────────────────────────────
\begin{center}
{\large\bfseries Слайд 2. Комбинирование и принятие решения}
\end{center}
\medskip
\hrule
\bigskip

\noindent\textbf{Шаг 4. Коэффициент конфликта и правило Демпстера}

\[
  K(m_1, m_2) =
    m_1\!\bigl(\{T\}\bigr) \cdot m_2\!\bigl(\{F\}\bigr)
    + m_1\!\bigl(\{F\}\bigr) \cdot m_2\!\bigl(\{T\}\bigr)
\]

\[
  (m_1 \oplus m_2)(A) =
  \frac{\displaystyle
    \sum_{\substack{B \cap C = A \\ A \neq \emptyset}}
      m_1(B) \cdot m_2(C)}{1 - K},
  \qquad K < \tau_K
\]

\[
  m^* = m'_{\mathrm{vac}} \oplus m'_{\mathrm{exp}} \oplus m'_{\mathrm{fc}},
  \qquad
  K = \max\!\bigl(K_1,\; K_2\bigr)
\]

\bigskip
\noindent\textbf{Шаг 5. Пигнистическая вероятность и дефицит}

\[
  \mathrm{BetP} = m^*\!\bigl(\{T\}\bigr) + \frac{m^*(\Theta)}{N_C},
  \qquad
  \Delta = \mathrm{BetP} - \sigma
\]

\noindent
$\sigma \in [0,1]$ --- текущая обеспеченность кластера по РПД/учебному плану.

\bigskip
\noindent\textbf{Шаг 6. Правила принятия решений}

\begin{align*}
  \Delta > \tau_\Delta
    \;\wedge\; K \le \tau_K
    \;\wedge\; m^*(\Theta) \le \tau_\Theta
  &\;\Longrightarrow\; \textbf{УСИЛИТЬ}\\[6pt]
  \Delta > \tau_\Delta
    \;\wedge\; \bigl(K > \tau_K
                     \;\vee\; m^*(\Theta) > \tau_\Theta\bigr)
  &\;\Longrightarrow\; \textbf{ВАРИАТИВНОСТЬ}\\[6pt]
  |\Delta| \le \tau_\Delta
  &\;\Longrightarrow\; \textbf{СТАБИЛИЗИРОВАТЬ}
\end{align*}

\smallskip
\noindent
$\tau_\Delta = 0{,}15$;\quad
$\tau_K = 0{,}40$;\quad
$\tau_\Theta = 0{,}15$;\quad
$N_C = 25$.

\medskip
\noindent\textit{Выход: рекомендация — УСИЛИТЬ / ВАРИАТИВНОСТЬ /
СТАБИЛИЗИРОВАТЬ / СОКРАТИТЬ.
При $K \ge \tau_K$ возможно переключение на правило Ягера
(конфликтная масса переходит на $\Theta$).}

\end{document}
